\chapter{Co-Estimation of States and Parameters}
\label{ch:co_estimation}

In advanced Battery Management Systems (BMS), state estimation—specifically State-of-Charge (SOC) and State-of-Health (SOH)—and internal equivalent circuit parameter tracking are deeply coupled. Internal parameters such as ohmic resistance $R_0$, polarization resistance $R_1$, polarization capacitance $C_1$, and total usable capacity $C_{\text{cap}}$ vary non-linearly as functions of SOC, cell temperature, and operational age. Concurrently, state estimators depend on precise parameter values to calculate open-circuit voltage $V_{oc}(\text{SOC})$ offsets and transient overvoltages.

This chapter explores joint and dual estimation architectures used to resolve state-parameter coupling. We analyze the physical origins of coupling, present simplified Joint EKF and Dual EKF algorithms, examine multi-time-scale estimation frameworks, and discuss practical implementation strategies for embedded microcontrollers.

% IMAGE PROMPT: High-level architectural flowchart of a Battery Management System (BMS) showing the mutual dependency between State Estimation (SOC, SOH) and Parameter Tracking (R0, R1, C1). Clean technical vector diagram with blue and green accent colors, illustrating closed-loop feedback between voltage measurements, state updates, and model adaptation.
% \begin{figure}[htbp]
%     \centering
%     \includegraphics[width=0.85\textwidth]{figures/ch7_coestimation_overview.png}
%     \caption{Overview of closed-loop mutual coupling between state estimation and parameter identification in a BMS.}
%     \label{fig:ch7_coestimation_overview}
% \end{figure}

\section{Coupling Between States (SOC/SOH) and Parameters}
\label{sec:coupling_states_parameters}

The state dynamic equation governing State-of-Charge progression under load current $I(t)$ is given by Coulomb counting:
\begin{equation}
\text{SOC}_k = \text{SOC}_{k-1} - \frac{\eta \cdot T_s}{C_{\text{cap}}} I_k
\label{eq:ch7_soc_coulomb}
\end{equation}
where $\eta$ is Coulombic efficiency, $T_s$ is the sampling interval, and $C_{\text{cap}}$ is the maximum usable capacity in Ampere-seconds.

Concurrently, the terminal voltage observation model depends explicitly on internal parameter values:
\begin{equation}
V_{t,k} = V_{oc}(\text{SOC}_k) - V_{1,k} - I_k R_0
\label{eq:ch7_vt_coupling}
\end{equation}

Equations~\eqref{eq:ch7_soc_coulomb} and~\eqref{eq:ch7_vt_coupling} establish a tight mutual feedback loop:
\begin{itemize}
    \item \textbf{State-to-Parameter Dependence:} Parameter estimation algorithms evaluate prediction residuals $e_k = V_{t,k} - \hat{V}_{t,k}$. If an SOC estimation error exists, errors propagate directly into $V_{oc}(\text{SOC})$, causing parameter identification algorithms to mistake SOC error for resistance growth or capacity fade.
    \item \textbf{Parameter-to-State Dependence:} Model-based state estimators rely on parameter values ($R_0, R_1, C_1$) to predict voltage drop and polarization dynamics. Inaccurate parameter values distort state updates, leading to persistent SOC tracking bias.
\end{itemize}

\subsection{Physical Insights and Identifiability Challenges}
In practical electric vehicle driving cycles, identifying both rapid state changes (such as polarization voltage transients) and slow parameter degradation (such as ohmic resistance growth over years) introduces numerical and observational challenges:
\begin{enumerate}
    \item \textbf{Time-Scale Separation:} Physical states like polarization voltage $V_1$ respond within seconds to current pulses, while capacity fade $C_{\text{cap}}$ and resistance increase $R_0$ evolve gradually over hundreds of charge-discharge cycles.
    \item \textbf{Flat OCV Region Sensitivity:} In lithium iron phosphate ($\text{LiFePO}_4$) cells, the slope $\frac{d V_{oc}}{d \text{SOC}}$ is near zero across 20\% to 80\% SOC. In this plateau region, small voltage measurement errors can cause large estimation drifts in both SOC and parameters if estimated simultaneously without regularization.
    \item \textbf{Excitation Requirements:} Parameters like $R_1$ and $C_1$ are observable only when the battery is subjected to dynamic current pulses. During long periods of constant-current charging or vehicle idling, parameter estimators lose excitation and must be suspended to prevent numerical divergence.
\end{enumerate}

\section{Joint Estimation Frameworks}
\label{sec:joint_estimation_frameworks}

Joint estimation frameworks tackle the coupled state and parameter tracking problem by concatenating physical states and unknown model parameters into a single augmented state vector evaluated by a unified filter.

\subsection{Augmented State Vector Formulation}
In the Joint Extended Kalman Filter (J-EKF) approach, the physical state vector $\mathbf{x}_k = [\text{SOC}_k, V_{1,k}]^T$ and parameter vector $\boldsymbol{\theta}_k = [R_{0,k}, R_{1,k}, C_{1,k}]^T$ are merged into an augmented state vector $\mathbf{x}_{a,k}$:
\begin{equation}
\mathbf{x}_{a,k} = \begin{bmatrix} \text{SOC}_k & V_{1,k} & R_{0,k} & R_{1,k} & C_{1,k} \end{bmatrix}^T
\label{eq:ch7_augmented_vector_def}
\end{equation}

Parameter dynamics are modeled as stationary random-walk processes driven by small fictitious noise $w_{\theta}$:
\begin{equation}
\boldsymbol{\theta}_k = \boldsymbol{\theta}_{k-1} + w_{\theta, k-1}
\label{eq:ch7_random_walk_params}
\end{equation}

The unified system step is written concisely as:
\begin{align}
\mathbf{x}_{a,k} &= f_a(\mathbf{x}_{a,k-1}, I_{k-1}) + \mathbf{w}_{a,k-1} \\
V_{t,k} &= g_a(\mathbf{x}_{a,k}, I_k) + v_k
\end{align}
where $\mathbf{w}_{a}$ represents process noise and $v_k$ is terminal voltage measurement noise.

% IMAGE PROMPT: Detailed structural diagram comparing Joint Estimation and Dual Estimation frameworks for Battery Management Systems. On the left, show a single unified Joint filter updating an augmented vector of states and parameters. On the right, show Dual estimation with two separate, coupled filters running at distinct time scales. Sleek vector layout.
% \begin{figure}[htbp]
%     \centering
%     \includegraphics[width=0.9\textwidth]{figures/ch7_joint_vs_dual.png}
%     \caption{Structural comparison between Joint Augmented Estimation (left) and Decoupled Dual Estimation (right).}
%     \label{fig:ch7_joint_vs_dual}
% \end{figure}

\subsection{Joint Extended Kalman Filter (J-EKF)}
The Joint EKF updates the augmented vector using standard Kalman prediction and correction steps:

\textbf{1. State Prediction:}
\begin{equation}
\hat{\mathbf{x}}_{a,k|k-1} = f_a(\hat{\mathbf{x}}_{a,k-1|k-1}, I_{k-1})
\end{equation}
\begin{equation}
\mathbf{P}_{a,k|k-1} = \mathbf{F}_{a,k-1} \mathbf{P}_{a,k-1|k-1} \mathbf{F}_{a,k-1}^T + \mathbf{Q}_a
\end{equation}

\textbf{2. Measurement Correction:}
\begin{equation}
e_k = V_{t,k} - g_a(\hat{\mathbf{x}}_{a,k|k-1}, I_k)
\end{equation}
\begin{equation}
\mathbf{K}_{a,k} = \mathbf{P}_{a,k|k-1} \mathbf{H}_{a,k}^T \left( \mathbf{H}_{a,k} \mathbf{P}_{a,k|k-1} \mathbf{H}_{a,k}^T + R_v \right)^{-1}
\end{equation}
\begin{equation}
\hat{\mathbf{x}}_{a,k|k} = \hat{\mathbf{x}}_{a,k|k-1} + \mathbf{K}_{a,k} e_k
\end{equation}
\begin{equation}
\mathbf{P}_{a,k|k} = (\mathbf{I} - \mathbf{K}_{a,k} \mathbf{H}_{a,k}) \mathbf{P}_{a,k|k-1}
\end{equation}

Here, $\mathbf{F}_a$ and $\mathbf{H}_a$ represent the linear derivative matrices of the system and measurement functions with respect to the augmented vector.

\begin{algorithm}[htbp]
\caption{Joint Extended Kalman Filter (J-EKF) Execution Steps}
\label{alg:joint_ekf}
\small
\begin{algorithmic}[1]
\State \textbf{Input:} Current $I_k$, Voltage $V_{t,k}$, noise covariances $\mathbf{Q}_a, R_v$
\State \textbf{Output:} Augmented state estimate $\hat{\mathbf{x}}_{a,k} = [\widehat{\text{SOC}}_k, \hat{V}_{1,k}, \hat{R}_{0,k}, \hat{R}_{1,k}, \hat{C}_{1,k}]^T$
\State \textbf{Step 1 (Predict):} Propagate states via Coulomb counting and RC dynamics; set parameter predictions equal to previous estimates.
\State \textbf{Step 2 (Covariance Predict):} Update augmented covariance matrix $\mathbf{P}_{a,k|k-1} = \mathbf{F}_a \mathbf{P}_{a,k-1|k-1} \mathbf{F}_a^T + \mathbf{Q}_a$.
\State \textbf{Step 3 (Innovation):} Compute voltage error $e_k = V_{t,k} - [V_{oc}(\widehat{\text{SOC}}_{k|k-1}) - \hat{V}_{1,k|k-1} - I_k \hat{R}_{0,k|k-1}]$.
\State \textbf{Step 4 (Correction):} Compute gain vector $\mathbf{K}_{a,k}$ and update all states and parameters simultaneously.
\end{algorithmic}
\end{algorithm}

\subsection{Practical Drawbacks of Joint Estimation}
While mathematically unified, joint estimation presents engineering challenges when implemented in production automotive microcontrollers:
\begin{itemize}
    \item \textbf{Scale Mismatch:} States like SOC lie in $[0, 1]$, while capacitances like $C_1$ range from $100\text{ F}$ to $5000\text{ F}$. Combining them in a single covariance matrix $\mathbf{P}_a$ leads to numerical ill-conditioning.
    \item \textbf{High Computational Overhead:} Inverting matrices for a 5-dimensional augmented vector requires significant floating-point operations per time step.
    \item \textbf{Fictitious Noise Tuning:} Setting the parameter process noise $\mathbf{Q}_{\theta}$ too high causes jittery parameter updates, while setting it too low freezes adaptation entirely.
\end{itemize}

\section{Dual Estimation Architecture}
\label{sec:dual_estimation_architecture}

To overcome the numerical scale mismatch of joint estimation, Dual Estimation decouples state estimation from parameter identification into two separate filters operating in an alternating, cooperative loop.

% IMAGE PROMPT: Block diagram of a Dual Extended Kalman Filter (DEKF) co-estimation system for lithium-ion batteries. Show two interconnected feedback loops: a fast State Filter running at 100 Hz estimating State of Charge (SOC) and polarization voltage, and a slow Parameter Filter running at 1 Hz updating internal resistance R0 and capacitance C1. Modern clean technical vector style with labeled signal flow arrows.
% \begin{figure}[htbp]
%     \centering
%     \includegraphics[width=0.85\textwidth]{figures/ch7_dekf_architecture.png}
%     \caption{Dual Extended Kalman Filter (DEKF) Co-Estimation Architecture with Multi-Timescale Decoupling.}
%     \label{fig:ch7_dekf_architecture}
% \end{figure}

\subsection{Multi-Time-Scale Decoupling Framework}
In the Dual Extended Kalman Filter (DEKF) framework:
\begin{enumerate}
    \item \textbf{State Filter (Fast Time-Scale):} Runs at a fast sampling rate (e.g., $100\text{ Hz}$ or $T_s = 10\text{ ms}$) to track rapidly changing states $\mathbf{x}_k = [\text{SOC}_k, V_{1,k}]^T$. It treats the current parameter estimates $\hat{\boldsymbol{\theta}}$ as known constants.
    \item \textbf{Parameter Filter (Slow Time-Scale):} Runs at a reduced sampling rate (e.g., $0.1\text{ Hz}$ to $1\text{ Hz}$ or $T_\theta = 1\text{ s}$) to track slowly evolving parameters $\boldsymbol{\theta}_k = [R_0, R_1, C_1]^T$. It uses the latest state estimates $\hat{\mathbf{x}}$ to evaluate prediction errors.
\end{enumerate}

\subsection{Dual EKF Execution Loop}
At each time step $k$, the dual filter executes the following sequence:

\textbf{1. State Filter Step:}
Using latest parameters $\hat{\boldsymbol{\theta}}_{k-1}$, update state estimate $\hat{\mathbf{x}}_k$:
\begin{equation}
\hat{\mathbf{x}}_{k|k} = \hat{\mathbf{x}}_{k|k-1} + \mathbf{K}_{x,k} \left( V_{t,k} - g(\hat{\mathbf{x}}_{k|k-1}, \hat{\boldsymbol{\theta}}_{k-1}, I_k) \right)
\end{equation}

\textbf{2. Parameter Filter Step (Executed every $M$ steps):}
Using updated states $\hat{\mathbf{x}}_k$, evaluate parameter innovation and update parameter vector $\hat{\boldsymbol{\theta}}_k$:
\begin{equation}
e_{\theta, k} = V_{t,k} - g(\hat{\mathbf{x}}_k, \hat{\boldsymbol{\theta}}_{k|k-1}, I_k)
\end{equation}
\begin{equation}
\hat{\boldsymbol{\theta}}_{k|k} = \hat{\boldsymbol{\theta}}_{k|k-1} + \mathbf{K}_{\theta, k} e_{\theta, k}
\end{equation}

By decoupling states and parameters:
\begin{itemize}
    \item The state filter matrix size is reduced to $2 \times 2$.
    \item The parameter filter matrix size is reduced to $3 \times 3$.
    \item Computational complexity per step drops dramatically from $\mathcal{O}(5^3)$ to $\mathcal{O}(2^3) + \frac{1}{M}\mathcal{O}(3^3)$.
\end{itemize}

\begin{algorithm}[htbp]
\caption{Multi-Time-Scale Dual EKF Framework}
\label{alg:dual_ekf_framework}
\small
\begin{algorithmic}[1]
\State \textbf{Initialize:} State vector $\hat{\mathbf{x}}_0$, Parameter vector $\hat{\boldsymbol{\theta}}_0$, Covariances $\mathbf{P}_{x,0}, \mathbf{P}_{\theta,0}$
\For{each fast sampling interval $k = 1, 2, \dots$}
    \State Measure current $I_k$ and terminal voltage $V_{t,k}$.
    \State Execute \textbf{State Filter}: Predict and correct $\hat{\mathbf{x}}_k = [\widehat{\text{SOC}}_k, \hat{V}_{1,k}]^T$ using fixed $\hat{\boldsymbol{\theta}}$.
    \If{$k \pmod M == 0$} \Comment{Slow time-scale parameter update}
        \State Execute \textbf{Parameter Filter}: Update $\hat{\boldsymbol{\theta}}_k = [\hat{R}_{0,k}, \hat{R}_{1,k}, \hat{C}_{1,k}]^T$ using latest $\hat{\mathbf{x}}_k$.
    \EndIf
\EndFor
\end{algorithmic}
\end{algorithm}

\section{Algorithmic Comparison and Engineering Considerations}
\label{sec:coest_comparison}

To assist BMS software architects in selecting the appropriate co-estimation approach, Table~\ref{tab:ch7_coest_comparison} compares Joint EKF, Dual EKF, and Hybrid recursive observer structures across critical engineering dimensions.

\begin{table}[htbp]
\centering
\caption{Comparative evaluation of co-estimation frameworks for embedded BMS.}
\label{tab:ch7_coest_comparison}
\small
\renewcommand{\arraystretch}{1.2}
\begin{tabularx}{\textwidth}{|p{3.0cm}|X|X|X|}
\hline
\textbf{Metric} & \textbf{Joint EKF (J-EKF)} & \textbf{Dual EKF (DEKF)} & \textbf{Hybrid EKF-RLS} \\
\hline
\textbf{Matrix Dimension} & Single $5 \times 5$ matrix & Decoupled $2 \times 2$ and $3 \times 3$ & $2 \times 2$ EKF + $3 \times 1$ RLS \\
\hline
\textbf{Execution Overhead} & High (Single fast rate) & Low (Multi-timescale) & Very Low \\
\hline
\textbf{Numerical Stability} & Prone to scale mismatch & Highly stable & High, if RLS reset used \\
\hline
\textbf{Tuning Complexity} & High ($\mathbf{Q}_a$ joint matrix) & Moderate (Decoupled noise) & Low (Independent settings) \\
\hline
\textbf{Best Use Case} & High-performance cloud BMS & Automotive edge microcontrollers & Cost-sensitive microcontrollers \\
\hline
\end{tabularx}
\end{table}

\subsection{Key Design Rules for Implementation}
When deploying co-estimation algorithms in production vehicles:
\begin{enumerate}
    \item \textbf{Excitation-Gated Updates:} Always freeze parameter filter updates when load current $|I_k| < I_{\text{threshold}}$ or during constant-current charge plateaus.
    \item \textbf{Parameter Bounding Box:} Hard-code physical saturation limits on parameter estimates (e.g., $0.5\text{ m}\Omega \le R_0 \le 50\text{ m}\Omega$) to prevent filter divergence caused by sensor noise spikes.
    \item \textbf{Thermal Compensation:} Parameter estimates should be normalized to a standard temperature (e.g., $25^\circ\text{C}$) using an Arrhenius correction factor before evaluating SOH capacity degradation.
\end{enumerate}

\section{Chapter Summary}
\label{sec:ch7_summary}
Co-estimation bridges the gap between state awareness and parameter identification. By replacing dense, high-dimensional joint formulations with multi-time-scale Dual EKF architectures, embedded Battery Management Systems achieve high estimation accuracy for SOC and SOH while maintaining low memory footprints and execution timing suitable for real-time automotive microcontrollers.
